\section{生成与理解为什么没有真正融合}

上一章解释静态图像的根本差异，本章看张祥雨在多模态系统里的实际挫折。他曾试图把图像生成、图像理解、文字生成和文字理解放进同一个系统：把图像 token 化，用语言式 Transformer 处理多模态，再接入 diffusion 等生成模块。结果是，理解模型越来越强，生成模型也越来越强，但放在一起没有出现 1+1 大于 2 的叠加效果。

最尖锐的证据是：把生成部分拿掉，理解能力并不受影响；把理解模型接到生成模型上，生成可控性仍然差。模型能识别某个视频违反物理常识，却无法阻止自己生成违反物理常识的视频。这说明两条分支还没有共享同一种内部思考过程。

\begin{figure}[H]
\centering
\includegraphics[width=0.94\textwidth]{generation-understanding-gap.png}
\caption{生成与理解未融合：两个分支都变强，但放在一起没有 1+1 大于 2。自制概念图，依据 00:29:10--00:38:45 对谈内容整理。}
\end{figure}

\begin{knowledgebox}{读图：系统连接不等于能力融合}
左边理解模型擅长看图、视频和语言对齐；右边生成模型擅长产生图像和视频。中间的 gap 不是工程接口，而是能力无法互相校正。真正的一体化，应当让理解约束生成，让生成过程反过来训练理解。
\end{knowledgebox}

\subsection{生成可控性的根本问题}

本节把“可控性差”拆开。图像和视频生成模型可以学习高概率视觉模式，但复杂场景需要遵守几何、物理、身份一致性、时间连续性和语义约束。语言提示并不总能把这些约束精确传入生成过程；理解模型也不一定能把“这不合理”转成生成模型内部的可执行修正。

\begin{knowledgebox}{术语消化：生成、理解、对齐、可控性}
\noindent\begin{tabularx}{\textwidth}{p{0.18\textwidth}p{0.36\textwidth}X}
\toprule
概念 & 含义 & 本期中的问题 \\
\midrule
生成 & 产生图像、视频或文本 & 可能逼真，但不一定可控。 \\
理解 & 识别对象、关系、动作和语义 & 可指出错误，但未必能指导生成。 \\
对齐 & 与人类意图和评价一致 & 图像对齐依赖语言、偏好和任务定义。 \\
可控性 & 按约束稳定生成目标结果 & 多物体、物理和时序约束最容易失败。 \\
\bottomrule
\end{tabularx}
\end{knowledgebox}

\begin{warningbox}{生成模型会“知道错”但仍然做错}
这不是表面矛盾。判断一个结果是否错，和在生成过程中一步步避免错误，是两种不同能力。前者像批改，后者像规划和执行。多模态系统缺的，正是连接批改和执行的思考链路。
\end{warningbox}

\subsection{本章小结}

多模态一体化的困难在于，理解和生成并不是接上线就能互相增强。它们需要共享可解释、可搜索、可纠错的中间过程；这正是后面 o1、CoT 和视觉 Long CoT 进入讨论的原因。

